\chapter{Run the Test from the Host}

The test is useful, but running it still requires opening a jail console and
remembering which services must already be available. We will put that
workflow behind one host-side command. Blackship will start the dependency
graph, the script will wait for the web edge, and the test jail will execute
the same Minitest file as before.

\section{Exercise 24A: Create the host runner}

On the FreeBSD host, become \texttt{root} and create
\texttt{/usr/local/sbin/rails-on-jails-test}:

\begin{lstlisting}[language=sh]
vi /usr/local/sbin/rails-on-jails-test
\end{lstlisting}

Enter:

\begin{lstlisting}[language=sh]
#!/bin/sh
set -eu

if [ "$(id -u)" -ne 0 ]; then
  echo "rails-on-jails-test must run as root" >&2
  exit 1
fi

blackship=/usr/local/bin/blackship
config_root=/usr/local/etc/bsdrunner
web_url=http://10.17.89.26/
test_file=/usr/local/tests/store_e2e_test.rb

set -- --config "${config_root}/blackship.toml"
for fragment in "${config_root}"/blackship.d/*.toml; do
  [ -f "${fragment}" ] || continue
  set -- "$@" --config "${fragment}"
done

"${blackship}" "$@" --non-interactive up test

attempt=1
while ! "${blackship}" "$@" --non-interactive exec test -- \
    /usr/bin/fetch -qo /dev/null "${web_url}" \
    >/dev/null 2>&1; do
  if [ "${attempt}" -ge 30 ]; then
    echo "web edge did not become ready within 60 seconds" >&2
    exit 1
  fi
  sleep 2
  attempt=$((attempt + 1))
done

exec "${blackship}" "$@" --non-interactive exec test -- \
  /usr/local/bin/ruby "${test_file}"
\end{lstlisting}

The \texttt{set --} loop builds a positional-argument list containing the
shared manifest followed by every per-jail manifest. Quoting \texttt{"\$@"}
passes those arguments to each Blackship invocation without flattening paths
into one string.

The readiness loop is separate from the test. Transient startup time can cause
retries, but once nginx returns a successful response, Minitest runs exactly
once. The final shell \texttt{exec} also makes the wrapper return Blackship's
test status unchanged, which is the behavior a CI job needs.

\section{Exercise 24B: Validate the script}

Make the runner executable and ask the shell to parse it without executing it:

\begin{lstlisting}[language=sh]
chmod 0555 /usr/local/sbin/rails-on-jails-test
sh -n /usr/local/sbin/rails-on-jails-test
\end{lstlisting}

No output from \texttt{sh -n} means the shell found no syntax errors. It does
not start jails or run the test.

\section{Exercise 24C: Run the complete workflow}

From a normal host account, use \texttt{mdo} to run the root-owned workflow:

\begin{lstlisting}[language=sh]
mdo /usr/local/sbin/rails-on-jails-test
\end{lstlisting}

If the current shell is already running as \texttt{root}, omit \texttt{mdo}.
Blackship may report that some or all jails are already running. That is
expected: \texttt{up test} starts only missing members of the declared
dependency graph. After the readiness check, expect one Minitest run, five
assertions, no failures, and no errors.

The command deliberately leaves the jails running. Automatically stopping the
graph would risk shutting down services that were active before the test began.
A CI environment that owns its entire host can add explicit teardown later.

\section{Exercise 24D: Restore a missing runner}

Use BSDRunner to stop only the \texttt{test} jail. It is the leaf of the
dependency graph, so this does not stop the application jails. Then invoke the
same host command again:

\begin{lstlisting}[language=sh]
mdo /usr/local/sbin/rails-on-jails-test
\end{lstlisting}

This time Blackship should report that the four application jails are already
running, start the \texttt{test} jail, and produce the same passing Minitest
result. No manual console or separate start action should be necessary.

\section{Checkpoint}

Stop after the command passes both with the complete graph already running and
with the test jail initially stopped. The manual console procedure is now a
repeatable command with a meaningful process exit status, ready to be called
by a scheduler or CI system.
